% LuaLaTeX + jlreq. TeX Live 2024 ships both jlreq and luatexja, and the
% HaranoAji fonts they default to, so this builds with no font installation.
\usepackage{luatexja}
\usepackage[haranoaji]{luatexja-preset}

\usepackage{amsmath}
\usepackage{amssymb}
\usepackage{graphicx}
\usepackage{booktabs}
\usepackage{longtable}
\usepackage{siunitx}
\usepackage{xcolor}
\usepackage[table]{xcolor}
\usepackage{caption}
\usepackage{subcaption}
\usepackage{float}
\usepackage[backend=bibtex,style=authoryear,natbib=true]{biblatex}
\usepackage[hidelinks]{hyperref}

\addbibresource{refs.bib}

% Figures live where the pipeline wrote them, so the report can never show a
% stale copy.
\graphicspath{{../out/figures/}{../out/figures/hokkaido/}{../out/figures/wakayama/}
  {../out/figures/hokkaido_week/}{../out/figures/wakayama_week/}
  {../out/figures/compare/}}

\sisetup{
  detect-all,
  per-mode = symbol,
  range-phrase = {--},
  range-units = single,
}

% Japanese unit names are written as \num{140} 分, not \SI{140}{分}. siunitx
% sets the unit argument in a Latin font, so a CJK unit is silently dropped and
% the reader is left with a bare number: the log says "Missing character: There
% is no 分 (U+5206) in font [lmroman10-regular]" and the PDF says "約 140".
% \DeclareSIUnit does not help -- the glyph is lost at typesetting, not at
% parsing -- and neither does wrapping it in \mbox or \text.

\captionsetup{
  font = small,
  labelfont = bf,
  justification = raggedright,
  singlelinecheck = false,
}

\definecolor{linkblue}{HTML}{2A78D6}
\hypersetup{
  colorlinks = true,
  linkcolor = linkblue,
  citecolor = linkblue,
  urlcolor = linkblue,
}

% Shorthands for the units that appear constantly.
\newcommand{\ms}{\si{\metre\per\second}}
\newcommand{\hpa}{\si{\hecto\pascal}}
\newcommand{\km}{\si{\kilo\metre}}

% A box for the "why this choice" asides, which are the point of this report.
\usepackage{tcolorbox}
\tcbuselibrary{breakable, skins}
\newtcolorbox{whybox}[1]{
  colback = linkblue!4,
  colframe = linkblue!40,
  boxrule = 0.4pt,
  left = 6pt, right = 6pt, top = 5pt, bottom = 5pt,
  title = {#1},
  fonttitle = \bfseries\small,
  breakable,
}

% Code blocks. listings rather than minted: minted shells out to pygmentize,
% which lives only in the conda environment, so a plain `make` outside it
% would fail. listings is pure TeX and always builds.
%
% Only the command names are keywords. Using listings' bash language would
% colour `cd` and `env` while leaving `conda` and `python` plain, which reads
% as random rather than as a rule.
\usepackage{listings}
\usepackage{needspace}
\definecolor{codebg}{HTML}{F7FAFD}
\definecolor{codecomment}{HTML}{6B7280}
\definecolor{codestring}{HTML}{9A4B00}
\lstset{
  basicstyle = \ttfamily\small,
  keywordstyle = \color{linkblue}\bfseries,
  % No \itshape: the Japanese comments would come out as slanted mincho.
  commentstyle = \color{codecomment},
  stringstyle = \color{codestring},
  morekeywords = {conda, python, pytest, make, cd, ls, cat},
  comment = [l]{\#},
  morestring = [b]",
  morestring = [b]',
  showstringspaces = false,
  extendedchars = true,
  columns = fullflexible,
  keepspaces = true,
  breaklines = true,
  breakindent = 2em,
  backgroundcolor = \color{codebg},
  frame = single,
  rulecolor = \color{linkblue!40},
  framerule = 0.4pt,
  framesep = 6pt,
  xleftmargin = 6pt,
  xrightmargin = 6pt,
  aboveskip = \medskipamount,
  belowskip = \medskipamount,
}
% A framed block split two lines before a page break looks like a mistake.
% Start one on a fresh page unless several lines fit.
\newcommand{\codeblock}{\Needspace*{6\baselineskip}}

% jlreq measures the width of the contents-page number box once, when the
% class loads and \thesection is still \arabic{section}. The appendix
% renumbers sections to "付録A", which no longer fits, so the number runs
% into the title ("付録 BCSV 表の一覧"). Widen the box from the appendix on.
\makeatletter
\newcommand{\appendixtocwidth}{%
  \renewcommand{\l@section}{\@dottedtocline{1}{0pt}{4.5em}}%
}
\makeatother
